\documentclass[10pt,a4paper]{article}
\usepackage{fontspec}
\usepackage[english,russian]{babel}
\babelfont{rm}{PT Serif}
\babelfont{sf}{PT Sans}
\babelfont{tt}[Scale=0.88]{PT Mono}
\usepackage[a4paper,top=20mm,bottom=22mm,left=20mm,right=20mm]{geometry}
\usepackage{amsmath}
\usepackage{booktabs}
\usepackage{array}
\usepackage{enumitem}
\usepackage{titlesec}
\usepackage{caption}
\usepackage{tikz}
\usetikzlibrary{arrows.meta}
\usepackage{xcolor}
\usepackage[hidelinks]{hyperref}
\usepackage{microtype}

\setlength{\parindent}{0pt}
\setlength{\parskip}{3pt plus 1pt}
\setlength{\emergencystretch}{4em}

\titleformat{\section}{\normalfont\fontsize{11.5}{14}\bfseries}{\thesection.}{0.5em}{}
\titlespacing*{\section}{0pt}{8pt plus 2pt minus 2pt}{3pt}
\titleformat{\subsection}{\normalfont\fontsize{10}{12}\bfseries}{\thesubsection.}{0.5em}{}
\titlespacing*{\subsection}{0pt}{5pt plus 2pt minus 1pt}{2pt}

\setlist[enumerate]{leftmargin=2.0em,itemsep=0.5pt,topsep=2pt,parsep=0pt}
\setlist[itemize]{leftmargin=1.5em,itemsep=0.5pt,topsep=2pt,parsep=0pt}
\captionsetup{font=small,labelfont=bf,labelsep=period,skip=3pt}
\newcolumntype{L}[1]{>{\raggedright\arraybackslash}p{#1}}
\setlength{\textfloatsep}{7pt plus 2pt minus 2pt}
\setlength{\floatsep}{8pt plus 2pt minus 2pt}

\pagestyle{plain}

\begin{document}

\begin{center}
{\fontsize{15}{18.5}\selectfont\bfseries Swarma: состояние реализации\par}
\vspace{4pt}
{\normalsize Отдельное приложение к работе ``Swarma: проверяемые границы доверия\\ для роёв автономных агентов'', версия 6.2.0\par}
\vspace{6pt}
{\normalsize Алексей Григорьев, Иван Опарин\\
Relux Works\\
\href{https://github.com/relux-works}{github.com/relux-works}\par}
\vspace{3pt}
{\normalsize Октябрь 2026\par}
\end{center}

\section{Назначение и шкала}

Работа описывает устройство Swarma и разбирает его гарантии как аргументы дизайна. Этот документ записывает, в каком состоянии находится каждая часть на октябрь 2026 года, и служит единственным местом статусов для версии 6.2 работы. Статус никогда не выше того, что говорит источник.

Шкала статусов такая. \emph{Работает}: часть есть в коде и используется в названных условиях. \emph{Спецификация}: часть описана в спецификации модуля, и это не свидетельство реализации. \emph{Решено}: владельцы приняли решение, кода нет. \emph{Предложено}: архитектурное предложение ожидает решения. \emph{Открыто}: решения нет.

\section{Части платформы}

\begin{table}[h]
\centering\small
\caption{Состояние реализации на октябрь 2026 года. Где состояние не сводится к одному слову, оно дано формулировкой.}
\label{tab:status}
\footnotesize
\begin{tabular}{L{6.3cm}L{2.5cm}L{6.9cm}}
\toprule
Часть & Статус & Граница \\
\midrule
Доска задач: Change Request, ревью, compare-and-set & работает & без сертификатов кворума \\
Контракт носителя: абстрактный носитель и профили протоколов & решено & спецификация контракта не опубликована \\
Профиль Matrix носителя: доставка, звонок и чтение почты & работает & используемый профиль; без ворот, грантов и восстановления \\
Профиль XMPP носителя & открыто & не определён и не квалифицирован: адреса, комнаты, вход и отзыв сессий, привязка конверта, удостоверение устройств для A1 \\
Подписанные записи и библиотека проверки & спецификация & библиотека без сети и ввода-вывода; потребители не включены \\
Ключник для Linux v0: ядро и регистрация & в коде, ждёт ревью & нет развёрнутой границы и квалификации \\
Подпись сообщений ключником & решено & кода нет \\
Контракт \texttt{swarma-msg/2} & предложено & на ревью содержания \\
Реестр владельца на узле доверия & решено & проект службы на ревью; узла нет \\
Протокол и сервер свидетеля отката & предложено & не реализованы и не развёрнуты \\
Хранитель контрольной точки вне узла & открыто & провайдер не выбран \\
Журнал ключевых событий в реестре до Swarma Chain & предложено & нужен рантайму раньше цепочки \\
Мост и ворота доверия & решено & кода и спецификации моста нет \\
Регистратор допуска к носителю & спецификация & сервер носителя не квалифицирован \\
Менеджер пользователей, запускатель, диспетчер & спецификация & станции не допускаются; квалификационные запуски на одной машине \\
Брокер учётных данных: контракт v0 & спецификация & аренда стороне агента в режиме с читаемым токеном \\
Брокер учётных данных: экспериментальное ядро & спецификация & код ядра службы и исполнителя внешнего токена на Linux есть; поставляемый адаптер владельца отказывает в выдаче, набор квалифицированных выпусков пуст, ядро ничего не выдаёт \\
Режим без секретов у агента: дизайн & решено & предложение одобрено владельцем; кода нет, ничего не квалифицировано \\
Шлюз учётных данных, служба действий, адаптеры сервисов & решено & кода нет; через шлюз только провайдеры моделей и хостинг кода без записи, любая запись только через службу действий \\
Контракт режима без секретов в спецификациях брокера, шлюза и адаптеров & открыто & служебные принципалы и аренды, право на использование через службу, вызов блокировки аккаунта; формат не опубликован \\
Заглушка в агентской среде с настроенным адресом провайдера & открыто & ни одна среда, адаптер и топология не квалифицированы \\
Растяжка утечки и блокировка аккаунта & решено & блокировка на весь аккаунт, снимает только владелец; подтверждённая утечка требует ротации у провайдера, возможная только блокирует; правило подтверждения и маршрут оповещения открыты \\
Сверка с выданными значениями вне шлюза, в том числе в мосте & открыто & граница сверки не одобрена; мост такой сверки не заявляет \\
Перехват TLS, обработка запросов в брокере, запись через шлюз & решено & исключены решением владельца, закрытые ворота \\
Личные подписки, другие классы сервисов через шлюз, донорские варианты B и C в режиме без секретов & открыто & закрытые ворота до квалификации на тестовых подписках, отдельного решения или отдельного шага с публикацией и квалификацией; первая поставка только локальная \\
Контракт изоляции исполнения & предложено & почти все ячейки матрицы \texttt{PENDING-VERIFY} \\
Проверка пакета памяти двумя вычислителями & работает & сам пакет решён; первой рабочей версии нет \\
Swarma Protocol, Swarma Chain & спроектированы, не построены & проверена ограниченная модель интеграции \\
\bottomrule
\end{tabular}
\end{table}

\section{На что опираются гарантии}

Гарантии работы получают силу не раньше своих механизмов. G1 опирается на ключник (в коде) и подпись сообщений (решено), G2 на контракт брокера v0 и невыбранный общий язык грантов, а в режиме без секретов у агента ещё на расширение контракта брокера для служебных аренд (открыто), G3 и G4 на реестр (решено) и свидетеля (предложено), G5 на мост (решено), G6 на регистратор (спецификация), G7 на контракт изоляции (предложено), G8 на доску (работает) и протокол (спроектирован), G9 на пакет памяти (решено). Режим без секретов у агента добавляет к G2, G3 и G7 зависимость от шлюза учётных данных и службы действий (решено, кода нет), от их проверки полномочия на каждом использовании и от квалификации сети шлюза; хранение учётных данных опирается на раздельные учётки этих служб и от root хоста не защищено. Большинство механизмов описано решением или спецификацией без кода.

\section{Открытые форматы и механизмы}

\begin{itemize}
\item Формат общего языка грантов не выбран; окончательный формат выбирает спецификация грантов.
\item Формат страниц снимка отзывов не определён; один полный снимок вмещает не больше 64 записей и 8192 байт.
\item Механизм завершения отзыва открыт.
\item Протокол свидетеля отката предложен, хранитель контрольной точки вне узла не выбран.
\item Сервер носителя не квалифицирован по девяти требованиям регистратора.
\item Версионированное расширение контракта брокера для режима без секретов у агента не опубликовано: служебные принципалы, служебные аренды, право на использование через службу и аутентифицированный вызов блокировки аккаунта.
\item Завершение отзыва у шлюза учётных данных и службы действий, то есть их долговечные подтверждения как потребителей отзыва, не опубликовано.
\item Граница сверки с выданными значениями вне шлюза и её владелец, правило подтверждённой утечки и маршрут оповещения владельца открыты.
\item Квалификация каждой агентской среды с заглушкой и настроенным адресом провайдера, каждого адаптера сервиса и каждой топологии открыта.
\end{itemize}

\section{Временные исключения и переходные политики}

В контракте брокера v0 действует одно временное исключение: ключ заявок вызывающего лежит в доме оркестратора, и его силу ограничивает грант. Запись git-профиля памяти ключник подписывает по запросу агента. Каналы брокера с читаемым токеном (окружение Claude, внешний токен Codex, ключ Muse) и собственный вход дома остаются явными исключениями из режима без секретов у агента по выбору владельца машины, и каждый квалифицируется отдельно. Профиль обновляющего Codex v0, где часть извлечения идёт под тем же UID, в гарантию режима без секретов не входит. Политика уровня A0, где отправителя-человека удостоверяет сервер носителя владельца, переходная. Профиль Matrix служит первым профилем носителя. Git-профиль памяти является первым профилем памяти. До запуска Swarma Chain журнал ключевых событий ведёт реестр владельца. План миграции долгоживущих корней на хэш-подписи предложен; ключи личности Swarma Chain спроектированы на SLH-DSA.

\section{Состояние свидетельств}

\begin{itemize}
\item Сквозного развёртывания, на котором можно измерить систему целиком, нет.
\item Сквозного прогона станции, где отзыв гранта или квалификации проходит через реестр, брокер, диспетчер и хост сессий и измеряется задержка до завершённого сужения или остановки, нет; граница 300~с и формула $\max(S+A,F)$ остаются свойствами дизайна.
\item Оценки производительности системы целиком и состязательной оценки против живого развёртывания нет.
\item Подписанной аттестации машины, которая называет каждую из 42 гарантий, нет, и почти все ячейки матрицы контракта изоляции остаются \texttt{PENDING-VERIFY}. Во втором привилегированном запуске допуск станций остался не настроен.
\item Сквозного пути режима без секретов у агента нет: шлюз, служба действий, заглушка и растяжка утечки не измерены, закрытие использования после отзыва не измерено. Синтетические проверки исполнителя, проверки контракта v0 и прошлые наблюдения входа и прокси эту квалификацию не заменяют.
\item Пакет памяти 0.3.0 на момент проверки вычислителем на Go был текущим черновиком; первой рабочей версии памяти нет.
\end{itemize}

\end{document}
